% PDF-to-TeX transcription by the assistant; not the original downloaded file.
\paragraph{Comparing two networks (Section 2.4.3).}
Given two sets of conductances $C$ and $\bar C$ on $G$, we say that $(G,\bar C)<(G,C)$ if $\bar C_{xy}<C_{xy}$ for all $xy$, or equivalently, if $\bar R_{xy}>R_{xy}$ for all $xy$. Assume that $(G,\bar C)<(G,C)$. Then by the Monotonicity Law, $\bar R_{\mathrm{eff}}\ge R_{\mathrm{eff}}$. Thus if random walk on $(G,\bar C)$ is transient, i.e., if $\bar R_{\mathrm{eff}}<\infty$, then random walk on $(G,C)$ is also transient. If random walk on $(G,C)$ is recurrent, i.e., if $R_{\mathrm{eff}}=\infty$, then random walk on $(G,\bar C)$ is also recurrent.
\begin{theorem}[Section 2.4.3]
If $(G,C)$ and $(G,\bar C)$ are networks, and if there exist constants $u,v$ with $0<u\le v<\infty$ such that
\[
uC_{xy}\le\bar C_{xy}\le vC_{xy}
\]
for all $x$ and $y$, then random walk on $(G,\bar C)$ is of the same type as random walk on $(G,C)$.
\end{theorem}
\begin{proof}
Let $U_{xy}=uC_{xy}$ and $V_{xy}=vC_{xy}$. Then $(G,U)\le(G,\bar C)\le(G,V)$. But the random walks for $(G,U)$ and $(G,V)$ are the same as random walk on $(G,C)$. Thus random walk for $(G,\bar C)$ is of the same type as random walk on $(G,C)$.
\end{proof}
\begin{corollary}[Section 2.4.3]
Let $(G,C)$ be a network. If for every edge $xy$ of $G$ we have $0<u<C_{xy}<v<\infty$ for some constants $u$ and $v$, then the random walk on $(G,C)$ has the same type as simple random walk on $G$.
\end{corollary}
\paragraph{The $k$-fuzz of a graph (Section 2.4.4).}
\begin{theorem}
Simple random walk on $G$ and on the $k$-fuzz $G_k$ of $G$ have the same type.
\end{theorem}
\begin{proof}
Let $\mathbf P$ be simple random walk on $G$. Define $\bar{\mathbf P}=(\mathbf P+\mathbf P^2+\ldots+\mathbf P^k)/k$. Then $\bar{\mathbf P}$ may be considered to be $\mathbf P$, watched at one of the first $k$ steps chosen at random, then at a time chosen at random from the next $k$ steps after this time, etc. Thinking of $\bar{\mathbf P}$ in this way, we see that $\mathbf P$ is in state $0$ at least once for every time $\bar{\mathbf P}$ in state $0$. Hence, if $\bar{\mathbf P}$ is recurrent so is $\mathbf P$. Assume now that $\bar{\mathbf P}$ is transient. Choose a finite set $S$ so that $0$ cannot be reached in $k$ steps from a point outside of $S$. Then, since the walk $\bar{\mathbf P}$ will be outside $S$ from some time on, the walk $\mathbf P$ cannot be at $0$ after this time, and $\mathbf P$ is also transient. Therefore, $\mathbf P$ and $\bar{\mathbf P}$ are of the same type.

Finally, we show that $\bar{\mathbf P}$ has the same type as simple random walk on $G_k$. Here it is important to remember our restriction that $G$ is of bounded degree, so that for some $E$ no vertex has degree $>E$. We know that $\mathbf P$ is reversible with $\mathbf w\mathbf P=\mathbf w$, where $w_x$ is the number of edges coming out of $x$. From its construction, $\bar{\mathbf P}$ is also reversible and $\mathbf w\bar{\mathbf P}=\mathbf w$. $\bar{\mathbf P}$ is the random walk on a network $(G_k,\bar C)$ with $\bar C_{xy}=w_x\bar P_{xy}$. If $\bar P_{xy}>0$, there is a path $x,x_1,x_2,\ldots,x_{m-1},y$ in $G$ from $x$ to $y$ of length $m\le k$. Then
\[
\bar P_{xy}\ge\frac1k\left(\frac1E\right)^m
\ge\frac1k\left(\frac1E\right)^k.
\]
Thus
\[
0<\frac1k\left(\frac1E\right)^k\le\bar P_{xy}\le1
\]
and
\[
0<\frac1k\left(\frac1E\right)^k\le\bar C_{xy}\le E.
\]
Therefore, by the theorem on the irrelevance of bounded twiddling proven in Section 2.4.3, $\bar{\mathbf P}$ and simple random walk on $G_k$ are of the same type. So $G$ and $G_k$ are of the same type.
\end{proof}
